% =====================================================================
%  Dokumentacja projektowa: HoneyWatch
%  Honeypot z analitycznym panelem monitorowania ataków sieciowych
%
%  Kompilacja:  pdflatex dokumentacja.tex  (dwukrotnie, dla spisu treści)
%  Wymaga: babel(polish), tikz, listings, hyperref, booktabs, geometry
% =====================================================================
\documentclass[12pt,a4paper]{report}

% ---------- Kodowanie i język ----------
\usepackage[utf8]{inputenc}
\usepackage[T1]{fontenc}
\usepackage[polish]{babel}
\usepackage{lmodern}

% ---------- Układ strony ----------
\usepackage[a4paper,left=2.5cm,right=2.5cm,top=2.5cm,bottom=2.5cm]{geometry}
\usepackage{setspace}
\onehalfspacing
\setlength{\parindent}{0pt}
\setlength{\parskip}{6pt}

% ---------- Grafika i kolory ----------
\usepackage{graphicx}
\usepackage{xcolor}
\definecolor{honey}{HTML}{F59E0B}
\definecolor{honeydark}{HTML}{B45309}
\definecolor{bgdark}{HTML}{1F2937}
\definecolor{codebg}{HTML}{F6F8FA}
\definecolor{codecomment}{HTML}{6A737D}
\definecolor{codekw}{HTML}{B45309}
\definecolor{codestr}{HTML}{0A7D2C}
\definecolor{linegray}{HTML}{9CA3AF}

% ---------- Tabele ----------
\usepackage{booktabs}
\usepackage{longtable}
\usepackage{array}
\usepackage{tabularx}
\usepackage{multirow}

% ---------- Listy ----------
\usepackage{enumitem}
\setlist{itemsep=2pt,topsep=3pt}

% ---------- Nagłówki/stopki ----------
\usepackage{fancyhdr}
\pagestyle{fancy}
\fancyhf{}
\fancyhead[L]{\small\itshape HoneyWatch}
\fancyhead[R]{\small\itshape Dokumentacja projektowa}
\fancyfoot[C]{\thepage}
\renewcommand{\headrulewidth}{0.4pt}
\renewcommand{\footrulewidth}{0pt}

% ---------- Podpisy rysunków/tabel ----------
\usepackage{caption}
\captionsetup{font=small,labelfont=bf,labelsep=period}

% ---------- TikZ ----------
\usepackage{tikz}
\usetikzlibrary{shapes.geometric,arrows.meta,positioning,fit,backgrounds,calc,shadows.blur}

% Style wspólne dla diagramów
\tikzset{
  block/.style={rectangle, rounded corners=3pt, draw=honeydark, thick,
                fill=honey!12, text width=3.2cm, align=center, minimum height=1.0cm,
                font=\small},
  store/.style={cylinder, shape border rotate=90, draw=honeydark, thick,
                fill=honey!20, aspect=0.3, minimum height=1.4cm, minimum width=2.6cm,
                align=center, font=\small},
  ext/.style={rectangle, draw=linegray, thick, fill=gray!8, text width=3.0cm,
              align=center, minimum height=0.9cm, font=\small},
  flow/.style={-{Stealth[length=3mm]}, thick, draw=honeydark},
  flowback/.style={-{Stealth[length=3mm]}, thick, draw=linegray, dashed},
  uclass/.style={rectangle, draw=honeydark, thick, fill=white, align=left,
                 font=\footnotesize, text width=4.4cm},
  actor/.style={font=\small, align=center},
  usecase/.style={ellipse, draw=honeydark, thick, fill=honey!12, align=center,
                  font=\footnotesize, minimum width=3.2cm, minimum height=1.0cm},
  lifeline/.style={rectangle, draw=honeydark, thick, fill=honey!12, align=center,
                   font=\footnotesize, minimum height=0.8cm, text width=2.4cm},
}

% ---------- Listingi kodu ----------
\usepackage{listings}
\lstset{
  basicstyle=\ttfamily\footnotesize,
  backgroundcolor=\color{codebg},
  commentstyle=\color{codecomment}\itshape,
  keywordstyle=\color{codekw}\bfseries,
  stringstyle=\color{codestr},
  numberstyle=\tiny\color{linegray},
  numbers=left,
  numbersep=8pt,
  frame=single,
  rulecolor=\color{linegray!60},
  breaklines=true,
  breakatwhitespace=false,
  showstringspaces=false,
  tabsize=2,
  captionpos=b,
  extendedchars=true,
  inputencoding=utf8,
  literate=%
    {ą}{{\k a}}1 {ć}{{\'c}}1 {ę}{{\k e}}1 {ł}{{\l}}1 {ń}{{\'n}}1
    {ó}{{\'o}}1 {ś}{{\'s}}1 {ź}{{\'z}}1 {ż}{{\.z}}1
    {Ą}{{\k A}}1 {Ć}{{\'C}}1 {Ę}{{\k E}}1 {Ł}{{\L}}1 {Ń}{{\'N}}1
    {Ó}{{\'O}}1 {Ś}{{\'S}}1 {Ź}{{\'Z}}1 {Ż}{{\.Z}}1
}
\lstdefinestyle{py}{language=Python}
\lstdefinestyle{sqlstyle}{language=SQL,morekeywords={AUTOINCREMENT,TEXT,INTEGER,DATETIME,UNIQUE}}
\lstdefinestyle{bashstyle}{language=bash}

% ---------- Hyperref (na końcu) ----------
\usepackage[hidelinks]{hyperref}
\hypersetup{
  pdftitle={Dokumentacja projektowa HoneyWatch},
  pdfauthor={Mateusz Rup, Jakub Gwizdala},
  colorlinks=true,
  linkcolor=honeydark,
  urlcolor=honeydark,
  citecolor=honeydark
}

% ---------- Makro: zastępczy zrzut ekranu ----------
% Aby wstawić rzeczywisty zrzut, zamień \zrzut{...}{...} na:
%   \includegraphics[width=\textwidth]{screenshots/plik.png}
\newcommand{\zrzut}[2]{%
  \begin{tikzpicture}
    \path[draw=linegray, fill=gray!6, rounded corners=4pt]
      (0,0) rectangle (\linewidth,#1)
      node[midway, align=center, text=linegray, font=\small, text width=0.85\linewidth]
      {[ Miejsce na zrzut ekranu ]\\[2pt]\itshape #2};
  \end{tikzpicture}%
}

\begin{document}

% =====================================================================
%  STRONA TYTUŁOWA
% =====================================================================
\begin{titlepage}
\centering
\thispagestyle{empty}

{\large Politechnika --- Wydział Informatyki\par}
\vspace{0.4cm}
{\large Kierunek: Informatyka\par}

\vfill

\begin{tikzpicture}
  \node[regular polygon, regular polygon sides=6, draw=honeydark, line width=2pt,
        fill=honey!25, minimum size=2.6cm] (hex) {};
  \node at (hex.center) {\Huge\bfseries\color{honeydark} HW};
\end{tikzpicture}

\vspace{0.8cm}

{\Huge\bfseries HoneyWatch\par}
\vspace{0.4cm}
{\Large Honeypot z analitycznym panelem\\ monitorowania ataków sieciowych\par}

\vspace{0.6cm}
{\large\itshape Dokumentacja techniczna i użytkowa projektu programistycznego\par}

\vfill

\begin{tabular}{rl}
\textbf{Autorzy:} & Mateusz Rup \\
                  & Jakub Gwizdała \\[6pt]
\textbf{Typ projektu:} & Projekt programistyczny \\[6pt]
\textbf{Repozytorium:} & \texttt{github.com/HoneyWatch/HoneyWatchApp} \\
\end{tabular}

\vspace{1.0cm}
{\large \today\par}

\vfill
\end{titlepage}

% =====================================================================
%  SPIS TREŚCI
% =====================================================================
\pagenumbering{roman}
\setcounter{page}{1}
\tableofcontents
\listoffigures
\cleardoublepage
\pagenumbering{arabic}
\setcounter{page}{1}

% =====================================================================
%  ROZDZIAŁ 1 — OPIS PROBLEMU
% =====================================================================
\chapter{Opis rozwiązywanego problemu}
\label{ch:problem}

\section{Cel projektu}

Celem projektu \textbf{HoneyWatch} było zaprojektowanie, zaimplementowanie oraz
wdrożenie kompletnego systemu typu \emph{honeypot} (pułapka sieciowa) wraz
z dedykowanym, analitycznym panelem (ang. \emph{dashboard}) służącym do
monitorowania i wizualizacji rzeczywistych ataków sieciowych prowadzonych przez
zautomatyzowane boty oraz ludzkich napastników.

W odróżnieniu od typowych projektów akademickich, w których dane są generowane
sztucznie lub pobierane z gotowych zbiorów, HoneyWatch zbiera dane
\emph{na żywo} --- system został wystawiony na publiczny adres IP serwera VPS
i przez cały okres działania rejestrował realne próby włamań. Dzięki temu projekt
łączy trzy obszary:

\begin{enumerate}
  \item \textbf{Infrastrukturę bezpieczeństwa} --- uruchomienie i konfigurację
        honeypotów emulujących usługi sieciowe (SSH, Telnet, FTP, HTTP, bazy
        danych i inne).
  \item \textbf{Inżynierię danych} --- parser logów przekształcający surowe
        zdarzenia w ustrukturyzowaną bazę danych, wzbogaconą o geolokalizację
        adresów IP.
  \item \textbf{Wizualizację i analizę} --- aplikację webową prezentującą dane
        w czasie zbliżonym do rzeczywistego: interaktywną mapę świata, wykresy,
        mapę cieplną aktywności oraz przeszukiwalną tabelę zdarzeń.
\end{enumerate}

Mówiąc krótko, chcieliśmy zbudować narzędzie, które odpowiada na pytanie:
\emph{,,kto, skąd, kiedy i w jaki sposób atakuje serwer wystawiony do
Internetu?''} --- i które prezentuje tę odpowiedź w sposób czytelny zarówno dla
specjalisty, jak i dla osoby niezwiązanej z bezpieczeństwem IT.

\section{Motywacja}

Każdy serwer udostępniony publicznie w Internecie jest niemal natychmiast
celem zautomatyzowanego skanowania i prób logowania. Dzieje się to w sposób
ciągły i masowy, jednak dla administratora zjawisko to pozostaje zwykle ukryte
w surowych, trudnych do interpretacji plikach logów. Motywacją do podjęcia
tematu była chęć:

\begin{itemize}
  \item \textbf{zobaczenia skali zjawiska na własne oczy} --- przekonania się,
        jak szybko i jak intensywnie ,,świeży'' serwer staje się celem ataków;
  \item \textbf{poznania technik napastników} --- jakie nazwy użytkowników
        i hasła są najczęściej próbowane, jakie usługi są skanowane, jakie
        polecenia wykonują napastnicy po uzyskaniu dostępu;
  \item \textbf{nauki praktycznego bezpieczeństwa} --- konfiguracji honeypotów,
        bezpiecznej izolacji usług oraz rejestrowania incydentów;
  \item \textbf{zbudowania użytecznego narzędzia} --- które przekształca ,,szum''
        z logów w wiedzę dającą się przedstawić wizualnie i wykorzystać do
        wyciągania wniosków obronnych.
\end{itemize}

Honeypot jest do tego celu narzędziem idealnym: jest to system, który
\emph{celowo} udaje podatną, wartościową maszynę, aby zwabić napastnika.
Ponieważ żaden legalny użytkownik nie ma powodu, by się do niego łączyć,
\textbf{każde} zarejestrowane połączenie można z dużym prawdopodobieństwem uznać
za wrogie. Eliminuje to fundamentalny problem klasycznych systemów wykrywania
włamań --- konieczność odróżniania ruchu legalnego od złośliwego.

\section{Czym jest honeypot}

\emph{Honeypot} to celowo wystawiona ,,przynęta'' --- usługa lub cały system
udający atrakcyjny cel (np. serwer SSH ze słabym hasłem), którego jedynym
zadaniem jest zwabienie napastnika i szczegółowe zarejestrowanie jego działań.
Wyróżnia się m.in. honeypoty:

\begin{itemize}
  \item \textbf{niskointeraktywne} --- emulujące jedynie wybrane elementy usługi
        (np. baner i ekran logowania), bezpieczne, lecz dostarczające mniej
        szczegółów;
  \item \textbf{wysokointeraktywne} --- udostępniające napastnikowi (pozornie)
        pełne środowisko, dające więcej danych kosztem większego ryzyka.
\end{itemize}

W projekcie wykorzystaliśmy dwa honeypoty o charakterze średnio-interaktywnym:
\textbf{Cowrie} (emulacja powłoki SSH/Telnet z rejestracją prób logowania
i wpisywanych poleceń) oraz \textbf{OpenCanary} (zestaw emulowanych usług na
wielu portach: HTTP, FTP, MySQL, RDP, VNC i innych). Dzięki ich połączeniu
rejestrujemy zarówno ataki na powłokę systemową, jak i skanowanie pozostałych
popularnych usług.

\section{Aktualność, oryginalność i użyteczność}

\textbf{Aktualność.} Bezpieczeństwo systemów wystawionych do Internetu jest
tematem nieustannie aktualnym. Skala zautomatyzowanych ataków rośnie wraz
z rozwojem botnetów, a analiza danych o atakach (\emph{threat intelligence})
stanowi jeden z filarów współczesnej obrony. Projekt dotyka więc realnego,
wciąż istotnego problemu.

\textbf{Oryginalność.} Choć istnieją gotowe, rozbudowane platformy do agregacji
danych z honeypotów (np. T-Pot, DShield), są to rozwiązania ciężkie, złożone
w konfiguracji i często ,,przeładowane'' funkcjami. Oryginalnym wkładem naszego
zespołu jest \textbf{lekki, samodzielnie napisany} potok przetwarzania danych
oraz dedykowany panel analityczny:

\begin{itemize}
  \item parser logów napisany od podstaw, bez zewnętrznych zależności
        (wyłącznie biblioteka standardowa Pythona), z mechanizmem
        \emph{przyrostowego} czytania logów i odporną na duplikaty zapisem;
  \item warstwa dostępu do danych odporna na różne formaty znaczników czasu
        i działająca identycznie na serwerze produkcyjnym oraz na komputerze
        deweloperskim;
  \item interfejs użytkownika zaprojektowany z naciskiem na czytelność ---
        ciemny motyw, interaktywna mapa, spójna kolorystyka i czytelne
        wskaźniki zmian.
\end{itemize}

\textbf{Użyteczność.} System dostarcza konkretną wartość: pozwala w kilka sekund
ocenić, z jakich krajów pochodzi najwięcej ataków, jakie poświadczenia są
najczęściej próbowane oraz w jakich godzinach aktywność jest największa. Tak
przetworzone dane mogą posłużyć do wzmocnienia rzeczywistych systemów (np.
blokowania określonych podsieci, wymuszania silnych haseł, ograniczania dostępu
do usług administracyjnych).

\section{Zakres i struktura dokumentu}

Niniejsza dokumentacja opisuje cały cykl życia projektu. W rozdziale
\ref{ch:tech} przedstawiono wykorzystane technologie. Rozdział \ref{ch:arch}
omawia architekturę systemu, model danych oraz kluczowe algorytmy i decyzje
projektowe, ilustrując je diagramami (blokowym, przypadków użycia, klas oraz
sekwencji). Rozdział \ref{ch:func} szczegółowo opisuje funkcjonalność z punktu
widzenia użytkownika, a rozdział \ref{ch:demo} prezentuje konkretne scenariusze
użycia wraz z wynikami. Rozdział \ref{ch:team} omawia organizację pracy zespołu,
a rozdział \ref{ch:summary} zawiera podsumowanie, wnioski, ograniczenia oraz
propozycje dalszego rozwoju. Dokument zamykają bibliografia oraz załącznik
z kodem źródłowym.

% =====================================================================
%  ROZDZIAŁ 2 — WYKORZYSTANE TECHNOLOGIE
% =====================================================================
\chapter{Wykorzystane technologie}
\label{ch:tech}

Zgodnie z zasadą, że dokumentacja powinna skupiać się na oryginalnym wkładzie
zespołu, w tym rozdziale ograniczamy się do wymienienia użytych narzędzi wraz
z krótkim uzasadnieniem wyboru. Szczegóły działania poszczególnych bibliotek
znajdują się w ich oficjalnej dokumentacji, do której odwołujemy się
w bibliografii.

\section{Zestawienie technologii}

\begin{longtable}{>{\raggedright\arraybackslash}p{3.4cm} >{\raggedright\arraybackslash}p{3.2cm} >{\raggedright\arraybackslash}p{7.2cm}}
\toprule
\textbf{Warstwa} & \textbf{Technologia} & \textbf{Rola w projekcie} \\
\midrule
\endfirsthead
\toprule
\textbf{Warstwa} & \textbf{Technologia} & \textbf{Rola w projekcie} \\
\midrule
\endhead
\bottomrule
\endfoot
Honeypoty & Cowrie & Emulacja powłoki SSH/Telnet; rejestracja prób logowania oraz poleceń. \\
          & OpenCanary & Emulacja wielu usług (HTTP, FTP, MySQL, RDP, VNC, SIP, MSSQL). \\
\midrule
Język & Python 3.10+ & Język parsera oraz aplikacji serwerowej. \\
\midrule
Parser & Biblioteka standardowa & \texttt{sqlite3}, \texttt{json}, \texttt{urllib}, \texttt{ipaddress}, \texttt{datetime} --- brak zależności zewnętrznych. \\
\midrule
Baza danych & SQLite & Lekka, plikowa baza relacyjna; jeden plik \texttt{honeywatch.db}. \\
\midrule
Backend WWW & Flask 3.x & Routing, renderowanie szablonów Jinja2, endpointy JSON. \\
\midrule
Frontend & HTML5 / CSS3 & Ciemny motyw, układ siatkowy (CSS Grid/Flexbox). \\
         & JavaScript (vanilla) & Pobieranie danych (\texttt{fetch}), renderowanie wizualizacji. \\
         & Chart.js 4.4 & Wykresy: liniowy, słupkowy, pierścieniowy. \\
         & Leaflet 1.9 & Interaktywna mapa świata z kafelkami CARTO Dark Matter. \\
\midrule
Geolokalizacja & ip-api.com & Wsadowe zapytania IP $\rightarrow$ kraj, szer./dł. geogr. \\
\midrule
Wdrożenie & Linux VPS & Serwer produkcyjny z publicznym adresem IP. \\
          & systemd & Usługa utrzymująca panel w działaniu. \\
          & cron & Cykliczne uruchamianie parsera (co minutę). \\
          & Bash & Skrypt instalacyjny \texttt{install-vps.sh}. \\
\midrule
Narzędzia & Git / GitHub & Kontrola wersji, współpraca, przegląd kodu (Pull Requests). \\
\end{longtable}

\section{Uzasadnienie kluczowych wyborów}

\begin{description}[leftmargin=0pt,style=nextline]
  \item[Cowrie i OpenCanary] Oba honeypoty są dojrzałe, dobrze udokumentowane
        i zapisują logi w formacie JSON, co znacznie upraszcza parsowanie.
        Cowrie specjalizuje się w SSH/Telnet, natomiast OpenCanary uzupełnia go
        o szeroki wachlarz innych usług --- razem dają pełny obraz aktywności.
  \item[Python + biblioteka standardowa] Parser celowo nie korzysta z bibliotek
        zewnętrznych (np. \texttt{requests}). Dzięki temu jego wdrożenie sprowadza
        się do skopiowania jednego pliku, bez instalowania zależności --- co jest
        istotne na minimalnym serwerze VPS.
  \item[SQLite] Dla skali tego projektu (pojedynczy serwer, dane rzędu tysięcy
        zdarzeń) plikowa baza SQLite jest w pełni wystarczająca, nie wymaga
        osobnego procesu serwera bazodanowego i ułatwia przenoszenie danych.
  \item[Flask] Minimalistyczny mikroframework, który nie narzuca struktury
        i pozwala szybko zbudować zarówno stronę HTML, jak i lekkie API JSON.
  \item[Chart.js i Leaflet] Sprawdzone, dojrzałe biblioteki frontendowe
        ładowane z sieci CDN, co odciąża serwer i nie wymaga procesu budowania
        (brak \texttt{npm}, brak kroku kompilacji).
\end{description}

% =====================================================================
%  ROZDZIAŁ 3 — ARCHITEKTURA I BUDOWA SYSTEMU
% =====================================================================
\chapter{Architektura i budowa systemu}
\label{ch:arch}

\section{Architektura ogólna}

System HoneyWatch ma architekturę \textbf{potoku przetwarzania danych}
(\emph{data pipeline}) złożoną z czterech głównych etapów: przechwytywania,
parsowania, składowania oraz prezentacji. Wszystkie komponenty działają na
jednym serwerze VPS, dzięki czemu dane nie muszą być przesyłane między
maszynami. Kluczową decyzją architektoniczną było to, że \textbf{panel czyta
plik bazy SQLite bezpośrednio}, w tym samym miejscu, w którym zapisuje go
parser --- eliminuje to potrzebę budowania pośredniej warstwy API między
parserem a panelem oraz ręcznego kopiowania bazy w środowisku produkcyjnym.

Schemat blokowy na rys.~\ref{fig:arch} przedstawia komponenty systemu oraz
kierunek przepływu danych.

\begin{figure}[ht]
\centering
\begin{tikzpicture}[node distance=1.0cm and 1.4cm]
  \node[ext] (att) {Napastnik / bot\\(Internet)};
  \node[block, below=of att] (hp) {Honeypoty\\\textbf{Cowrie} + \textbf{OpenCanary}};
  \node[store, below=of hp] (logs) {Pliki logów\\(JSON)};
  \node[block, below=of logs] (parser) {Parser\\\texttt{parser.py}\\(cron, co minutę)};
  \node[store, below=of parser] (db) {Baza danych\\\texttt{honeywatch.db}\\(SQLite)};
  \node[block, right=2.6cm of db] (flask) {Aplikacja Flask\\\texttt{dashboard}\\(systemd)};
  \node[ext, right=2.6cm of parser] (geo) {API geolokalizacji\\ip-api.com};
  \node[block, above=of flask] (geomod) {Moduł geoip\\+ cache JSON};
  \node[ext, above=of geomod] (browser) {Przeglądarka\\(klient)};

  \draw[flow] (att) -- node[right,font=\scriptsize] {ataki} (hp);
  \draw[flow] (hp) -- node[right,font=\scriptsize] {zapis} (logs);
  \draw[flow] (logs) -- node[right,font=\scriptsize] {odczyt przyrostowy} (parser);
  \draw[flow] (parser) -- node[right,font=\scriptsize] {INSERT OR IGNORE} (db);
  \draw[flow] (db) -- node[above,font=\scriptsize] {SELECT} (flask);
  \draw[flow] (flask) -- (geomod);
  \draw[flow] (geomod) -- node[left,font=\scriptsize] {nieznane IP} (geo);
  \draw[flow] (geomod) -- node[right,font=\scriptsize] {JSON / HTML} (browser);
\end{tikzpicture}
\caption{Architektura ogólna systemu HoneyWatch --- potok przetwarzania danych
od napastnika do przeglądarki.}
\label{fig:arch}
\end{figure}

\section{Przepływ danych}

Pełny cykl życia pojedynczego zdarzenia przebiega następująco:

\begin{enumerate}
  \item \textbf{Przechwycenie.} Napastnik (najczęściej zautomatyzowany bot)
        łączy się z serwerem na jednym z monitorowanych portów. Honeypot
        odpowiada, udając prawdziwą usługę, i zapisuje zdarzenie do pliku logu
        w formacie JSON (jedno zdarzenie = jedna linia).
  \item \textbf{Parsowanie.} Uruchamiany cyklicznie przez \texttt{cron} parser
        czyta \emph{tylko nowe} fragmenty logów (na podstawie zapamiętanego
        przesunięcia bajtowego), mapuje każde zdarzenie na model danych
        i zapisuje je do bazy z pominięciem duplikatów.
  \item \textbf{Składowanie.} Dane trafiają do relacyjnej bazy SQLite
        w czterech tabelach (zob. \ref{sec:schema}).
  \item \textbf{Wzbogacanie.} Przy obsłudze żądania panel pobiera unikalne
        adresy IP i ustala ich położenie geograficzne; wyniki są trwale
        zapisywane w pamięci podręcznej, aby ograniczyć liczbę zapytań do
        zewnętrznego API.
  \item \textbf{Prezentacja.} Aplikacja Flask agreguje dane (SQL) i zwraca je
        jako HTML (pierwsze załadowanie strony) oraz JSON (dynamiczne wykresy
        i mapa). Przeglądarka renderuje wizualizacje.
\end{enumerate}

\section{Diagram przypadków użycia}

Z perspektywy użytkownika system udostępnia zestaw funkcji analitycznych
skupionych wokół jednego panelu. Aktorami są \textbf{Analityk} (główny
użytkownik panelu) oraz dwa aktory systemowe: \textbf{Napastnik} (generujący
dane) i zewnętrzne \textbf{API geolokalizacji}. Rysunek \ref{fig:usecase}
przedstawia diagram przypadków użycia.

\begin{figure}[ht]
\centering
\begin{tikzpicture}[node distance=0.55cm]
  % Aktor: analityk
  \node[actor] (analyst) at (0,0) {%
    \begin{tikzpicture}[scale=0.5]
      \draw[thick] (0,0) circle (0.5);
      \draw[thick] (0,-0.5) -- (0,-1.8);
      \draw[thick] (-0.8,-1.0) -- (0.8,-1.0);
      \draw[thick] (0,-1.8) -- (-0.6,-2.6);
      \draw[thick] (0,-1.8) -- (0.6,-2.6);
    \end{tikzpicture}\\ \textbf{Analityk}};

  % Aktor: napastnik
  \node[actor] (attacker) at (13,1.5) {%
    \begin{tikzpicture}[scale=0.5]
      \draw[thick] (0,0) circle (0.5);
      \draw[thick] (0,-0.5) -- (0,-1.8);
      \draw[thick] (-0.8,-1.0) -- (0.8,-1.0);
      \draw[thick] (0,-1.8) -- (-0.6,-2.6);
      \draw[thick] (0,-1.8) -- (0.6,-2.6);
    \end{tikzpicture}\\ \textbf{Napastnik}};

  % Aktor: API
  \node[actor] (api) at (13,-4.5) {%
    \begin{tikzpicture}[scale=0.5]
      \draw[thick] (0,0) circle (0.5);
      \draw[thick] (0,-0.5) -- (0,-1.8);
      \draw[thick] (-0.8,-1.0) -- (0.8,-1.0);
      \draw[thick] (0,-1.8) -- (-0.6,-2.6);
      \draw[thick] (0,-1.8) -- (0.6,-2.6);
    \end{tikzpicture}\\ \textbf{API geolokalizacji}};

  % System boundary
  \begin{scope}[on background layer]
    \node[draw=honeydark, thick, rounded corners, fill=honey!4,
          minimum width=7.2cm, minimum height=8.6cm, label={[font=\small\bfseries]north:System HoneyWatch}]
          (sys) at (6.5,-1.5) {};
  \end{scope}

  \node[usecase] (uc1) at (6.5,1.3) {Przegląd wskaźników\\KPI};
  \node[usecase] (uc2) at (6.5,0.0) {Analiza mapy\\ataków};
  \node[usecase] (uc3) at (6.5,-1.3) {Przegląd wykresów\\i statystyk};
  \node[usecase] (uc4) at (6.5,-2.6) {Analiza\\poświadczeń};
  \node[usecase] (uc5) at (6.5,-3.9) {Przeszukiwanie\\logów};
  \node[usecase] (uc6) at (6.5,-5.2) {Zmiana zakresu\\czasu};

  \foreach \u in {uc1,uc2,uc3,uc4,uc5,uc6} {
    \draw[thick] (analyst) -- (\u);
  }
  \draw[thick] (attacker) -- node[above,sloped,font=\scriptsize]{generuje dane} (uc2);
  \draw[thick,dashed,-{Stealth}] (uc2) -- node[below,sloped,font=\scriptsize]{$\langle\langle$include$\rangle\rangle$} (api);
\end{tikzpicture}
\caption{Diagram przypadków użycia systemu HoneyWatch.}
\label{fig:usecase}
\end{figure}

\section{Model danych i schemat bazy}
\label{sec:schema}

Baza danych SQLite składa się z czterech tabel. Centralną encją jest tabela
\texttt{attacks}, do której odwołują się tabele szczegółowe \texttt{commands}
(polecenia powłoki) oraz \texttt{http\_attacks} (szczegóły żądań HTTP). Tabela
\texttt{parser\_state} przechowuje stan parsera (przesunięcia bajtowe).
Rysunek \ref{fig:erd} przedstawia uproszczony diagram związków encji.

\begin{figure}[ht]
\centering
\begin{tikzpicture}[node distance=1.2cm and 1.8cm]
  \node[uclass] (attacks) {%
    \textbf{attacks}\\\hrule\vspace{2pt}
    \underline{id} : INTEGER PK\\
    timestamp : TEXT\\
    src\_ip : TEXT\\
    src\_port : INTEGER\\
    username : TEXT\\
    password : TEXT\\
    success : INTEGER\\
    event\_type : TEXT\\
    source : TEXT};

  \node[uclass, right=of attacks] (commands) {%
    \textbf{commands}\\\hrule\vspace{2pt}
    \underline{id} : INTEGER PK\\
    attack\_id : INTEGER FK\\
    timestamp : TEXT\\
    src\_ip : TEXT\\
    command : TEXT};

  \node[uclass, below=of commands] (http) {%
    \textbf{http\_attacks}\\\hrule\vspace{2pt}
    \underline{id} : INTEGER PK\\
    attack\_id : INTEGER FK\\
    timestamp : TEXT\\
    src\_ip : TEXT\\
    method : TEXT\\
    path : TEXT\\
    attack\_type : TEXT};

  \node[uclass, below=of attacks] (state) {%
    \textbf{parser\_state}\\\hrule\vspace{2pt}
    \underline{log\_path} : TEXT PK\\
    byte\_offset : INTEGER};

  \draw[flow] (commands) -- node[above,font=\scriptsize]{1..*} (attacks);
  \draw[flow] (http) -- node[below,font=\scriptsize]{1..*} (attacks);
\end{tikzpicture}
\caption{Diagram związków encji bazy danych \texttt{honeywatch.db}.}
\label{fig:erd}
\end{figure}

Definicję tabel realizuje funkcja \texttt{init\_db()} (listing \ref{lst:schema}).

\begin{lstlisting}[style=sqlstyle,caption={Definicja schematu bazy danych (fragment).},label={lst:schema}]
CREATE TABLE IF NOT EXISTS attacks (
    id          INTEGER PRIMARY KEY AUTOINCREMENT,
    timestamp   TEXT,
    src_ip      TEXT,
    src_port    INTEGER,
    username    TEXT,
    password    TEXT,
    success     INTEGER,   -- 0 = nieudane, 1 = udane, NULL = nie dotyczy
    event_type  TEXT,      -- connect | login_failed | login_success | http_probe | ftp | ...
    source      TEXT       -- cowrie | opencanary
);

CREATE TABLE IF NOT EXISTS parser_state (
    log_path    TEXT PRIMARY KEY,
    byte_offset INTEGER NOT NULL DEFAULT 0
);
\end{lstlisting}

Znaczenie najważniejszych kolumn tabeli \texttt{attacks} przedstawia
tabela \ref{tab:cols}.

\begin{table}[ht]
\centering
\small
\begin{tabularx}{\textwidth}{>{\ttfamily}l X}
\toprule
\textnormal{\textbf{Kolumna}} & \textbf{Znaczenie} \\
\midrule
timestamp & Znacznik czasu zdarzenia (ISO-8601 lub epoka uniksowa). \\
src\_ip & Adres IP źródła ataku --- klucz do geolokalizacji. \\
src\_port & Port źródłowy połączenia. \\
username / password & Próbowane poświadczenia (dla zdarzeń logowania). \\
success & Wynik próby logowania (0/1/NULL). \\
event\_type & Typ zdarzenia, np. \texttt{connect}, \texttt{login\_failed}, \texttt{http\_probe}, \texttt{ftp}, \texttt{mysql}, \texttt{rdp}. \\
source & Honeypot będący źródłem zdarzenia (\texttt{cowrie}/\texttt{opencanary}). \\
\bottomrule
\end{tabularx}
\caption{Opis kolumn centralnej tabeli \texttt{attacks}.}
\label{tab:cols}
\end{table}

\section{Diagram klas / komponentów}

Logika aplikacji jest podzielona na spójne moduły o jasno określonej
odpowiedzialności. Rysunek \ref{fig:class} przedstawia uproszczony diagram
najważniejszych modułów i ich powiązań.

\begin{figure}[ht]
\centering
\begin{tikzpicture}[node distance=1.0cm and 1.0cm]
  \node[uclass, text width=4.6cm] (app) {%
    \textbf{dashboard.app} \emph{(Flask)}\\\hrule\vspace{2pt}
    + overview()\\
    + api\_summary()\\
    + api\_geo() / api\_timeline()\\
    + api\_event\_types()\\
    + api\_top\_usernames/passwords()\\
    + api\_heatmap() / api\_recent()\\
    + api\_logs()\\
    -- \_range()};

  \node[uclass, text width=4.6cm, right=of app] (db) {%
    \textbf{dashboard.db} \emph{(dostęp do danych)}\\\hrule\vspace{2pt}
    + get\_summary()\\
    + get\_top\_countries()\\
    + get\_geo() / get\_timeline()\\
    + get\_event\_types()\\
    + get\_heatmap() / get\_recent()\\
    + get\_logs()\\
    -- \_window\_sql() / \_ts\_mod()\\
    -- \_country\_aggregate()};

  \node[uclass, text width=4.6cm, below=of db] (geoip) {%
    \textbf{dashboard.geoip}\\\hrule\vspace{2pt}
    + locate(ips)\\
    + flag(code)\\
    -- \_fetch\_batch()\\
    -- \_load\_cache() / \_save\_cache()\\
    -- \_is\_public()};

  \node[uclass, text width=4.6cm, below=of app] (iso) {%
    \textbf{dashboard.iso\_codes}\\\hrule\vspace{2pt}
    + to\_iso3(alpha2)};

  \node[uclass, text width=4.6cm, left=of iso] (parser) {%
    \textbf{deploy.parser}\\\hrule\vspace{2pt}
    + init\_db()\\
    + parse\_cowrie()\\
    + parse\_opencanary()\\
    -- \_iter\_new\_lines()\\
    -- \_get/\_set\_offset()};

  \draw[flow] (app) -- node[above,font=\scriptsize]{wywołuje} (db);
  \draw[flow] (db) -- node[right,font=\scriptsize]{używa} (geoip);
  \draw[flow] (db) -- node[left,font=\scriptsize]{ISO3} (iso);
  \draw[flow] (parser) -- node[below,sloped,font=\scriptsize]{zapisuje DB} (db);
\end{tikzpicture}
\caption{Diagram głównych modułów systemu i ich zależności.}
\label{fig:class}
\end{figure}

\subsection{Opis modułów}

\begin{description}[leftmargin=0pt,style=nextline]
  \item[\texttt{deploy/parser.py}] Niezależny skrypt zbierający dane.
        Inicjalizuje bazę, czyta przyrostowo logi Cowrie i OpenCanary, mapuje
        zdarzenia na wiersze tabel i zapisuje je z deduplikacją. Uruchamiany
        cyklicznie przez \texttt{cron}.
  \item[\texttt{dashboard/app.py}] Cienka warstwa kontrolera Flask. Definiuje
        jedną stronę HTML oraz dziesięć endpointów JSON; całą logikę danych
        deleguje do modułu \texttt{db}.
  \item[\texttt{dashboard/db.py}] Serce analityczne aplikacji. Zawiera całą
        logikę SQL i agregacji: wskaźniki KPI, szeregi czasowe, mapę cieplną,
        rankingi, paginowane logi. Obsługuje zakresy czasu i automatyczną
        detekcję formatu znacznika czasu.
  \item[\texttt{dashboard/geoip.py}] Geolokalizacja adresów IP z trwałą pamięcią
        podręczną (\texttt{data/geo\_cache.json}); zapytania wsadowe do ip-api.com,
        pomijanie adresów prywatnych, generowanie emoji flag.
  \item[\texttt{dashboard/iso\_codes.py}] Tablica odwzorowań kodów ISO 3166-1
        alpha-2 $\rightarrow$ alpha-3, wymaganych przez mapę GeoJSON.
  \item[\texttt{dashboard/templates/}, \texttt{static/}] Warstwa prezentacji:
        szablony Jinja2 oraz pliki CSS i JavaScript renderujące mapę, wykresy
        i tabele.
\end{description}

\section{Diagram sekwencji}

Rysunek \ref{fig:seq} ilustruje przebieg typowej interakcji: załadowanie panelu
i pobranie danych do mapy ataków.

\begin{figure}[ht]
\centering
\begin{tikzpicture}
  % Lifelines
  \node[lifeline] (br) at (0,0) {Przeglądarka};
  \node[lifeline] (fl) at (4,0) {Flask\\\texttt{app.py}};
  \node[lifeline] (db) at (8,0) {Warstwa DB\\\texttt{db.py}};
  \node[lifeline] (gx) at (12,0) {geoip\\+ cache};

  \foreach \n in {br,fl,db,gx} {
    \draw[dashed, linegray] (\n.south) -- ($(\n.south)+(0,-9)$);
  }

  \def\y{-1}
  \draw[flow] ($(br.south)+(0,\y)$) -- node[above,font=\scriptsize]{GET /} ($(fl.south)+(0,\y)$);
  \def\y{-2}
  \draw[flow] ($(fl.south)+(0,\y)$) -- node[above,font=\scriptsize]{get\_summary()} ($(db.south)+(0,\y)$);
  \def\y{-2.8}
  \draw[flowback] ($(db.south)+(0,\y)$) -- node[above,font=\scriptsize]{HTML (KPI)} ($(fl.south)+(0,\y)$);
  \def\y{-3.6}
  \draw[flowback] ($(fl.south)+(0,\y)$) -- node[above,font=\scriptsize]{strona HTML} ($(br.south)+(0,\y)$);
  \def\y{-4.6}
  \draw[flow] ($(br.south)+(0,\y)$) -- node[above,font=\scriptsize]{GET /api/geo} ($(fl.south)+(0,\y)$);
  \def\y{-5.4}
  \draw[flow] ($(fl.south)+(0,\y)$) -- node[above,font=\scriptsize]{get\_geo()} ($(db.south)+(0,\y)$);
  \def\y{-6.2}
  \draw[flow] ($(db.south)+(0,\y)$) -- node[above,font=\scriptsize]{locate(ips)} ($(gx.south)+(0,\y)$);
  \def\y{-7.0}
  \draw[flowback] ($(gx.south)+(0,\y)$) -- node[above,font=\scriptsize]{kraj, lat/lng} ($(db.south)+(0,\y)$);
  \def\y{-7.8}
  \draw[flowback] ($(db.south)+(0,\y)$) -- node[above,font=\scriptsize]{agregaty} ($(fl.south)+(0,\y)$);
  \def\y{-8.6}
  \draw[flowback] ($(fl.south)+(0,\y)$) -- node[above,font=\scriptsize]{JSON} ($(br.south)+(0,\y)$);
\end{tikzpicture}
\caption{Diagram sekwencji: załadowanie panelu i pobranie danych mapy.}
\label{fig:seq}
\end{figure}

\section{Kluczowe algorytmy i decyzje projektowe}

W tej sekcji opisujemy oryginalne rozwiązania techniczne, które stanowią
najważniejszy wkład inżynierski zespołu.

\subsection{Przyrostowe czytanie logów}

Pliki logów rosną nieustannie, a parser uruchamiany jest co minutę. Czytanie
całego pliku za każdym razem byłoby marnotrawstwem. Dlatego parser zapamiętuje
w tabeli \texttt{parser\_state} \emph{przesunięcie bajtowe} (\texttt{byte\_offset})
do którego doczytał, i przy kolejnym uruchomieniu zaczyna od tego miejsca.
Mechanizm uwzględnia również \emph{rotację logów}: jeśli plik jest krótszy niż
zapisane przesunięcie, oznacza to, że został on obrócony, więc odczyt zaczyna
się od początku. Dodatkowo przetwarzane są tylko pełne linie (do ostatniego
znaku nowej linii), aby uniknąć parsowania linii zapisywanej właśnie przez
honeypot. Fragment tej logiki przedstawia listing \ref{lst:incr}.

\begin{lstlisting}[style=py,caption={Przyrostowe czytanie nowych linii logu (fragment).},label={lst:incr}]
def _iter_new_lines(c, log_path):
    offset = _get_offset(c, log_path)
    file_size = os.path.getsize(log_path)
    if file_size < offset:        # log zostal obrocony -> czytaj od poczatku
        offset = 0
    with open(log_path, "rb") as f:
        f.seek(offset)
        data = f.read()
    last_nl = data.rfind(b"\n")   # przetwarzaj tylko pelne linie
    if last_nl == -1:
        return
    chunk = data[:last_nl + 1]
    _set_offset(c, log_path, offset + len(chunk))
    for raw in chunk.split(b"\n"):
        line = raw.decode("utf-8", errors="replace").strip()
        if line:
            yield line
\end{lstlisting}

\subsection{Idempotentna deduplikacja zapisu}

Aby ten sam atak nie został policzony wielokrotnie (np. przy ponownym
przetworzeniu logu), na tabelach utworzono \emph{unikalne indeksy} obejmujące
zestaw kolumn jednoznacznie identyfikujących zdarzenie, a zapis wykonywany jest
poleceniem \texttt{INSERT OR IGNORE}. Jeśli w momencie tworzenia indeksu w bazie
istnieją już duplikaty (z wcześniejszych wersji parsera), funkcja
\texttt{\_dedupe\_existing()} usuwa nadmiarowe wiersze (zachowując rekord
o najniższym \texttt{id}) i ponawia próbę. Dzięki temu cały proces parsowania
jest \textbf{idempotentny} --- wielokrotne uruchomienie nie zniekształca danych.

\subsection{Trwała pamięć podręczna geolokalizacji}

Darmowy plan ip-api.com jest limitowany liczbą zapytań. Moduł \texttt{geoip}
ogranicza ruch sieciowy do absolutnego minimum:

\begin{itemize}
  \item każde unikalne IP jest odpytywane \emph{co najwyżej raz}, a wynik
        zapisywany w pliku \texttt{geo\_cache.json};
  \item adresy prywatne i zarezerwowane są rozpoznawane lokalnie
        (\texttt{ipaddress.is\_global}) i nigdy nie są wysyłane do API;
  \item zapytania są \emph{wsadowe} (do 100 adresów w jednym żądaniu);
  \item adresy, których nie udało się rozwiązać, są zapamiętywane jako
        ,,rozwiązane, lecz nieznane'' (\texttt{None}), aby nie próbować ich
        ponownie;
  \item wszystkie błędy sieciowe są wyciszane --- brak łączności z API powoduje
        jedynie brak flag, a nie awarię panelu.
\end{itemize}

\subsection{Niezależność od formatu znacznika czasu}

Różne honeypoty mogą zapisywać czas jako tekst ISO-8601 albo jako liczbę
(epoka uniksowa). Funkcja \texttt{\_ts\_mod()} próbkuje jeden znacznik czasu
z bazy i ustala, czy jest on liczbą; na tej podstawie dobiera modyfikator
(\texttt{'unixepoch'} lub jego brak), który jest następnie wstawiany do każdego
wywołania \texttt{datetime()}/\texttt{strftime()} w zapytaniach SQL. Dzięki temu
ta sama logika działa poprawnie niezależnie od formatu danych.

\subsection{Odporność na nieaktualną kopię bazy}

Funkcja \texttt{\_anchor()} wyznacza punkt odniesienia ,,teraz'' jako maksimum
z rzeczywistego czasu bieżącego oraz najnowszego znacznika czasu w bazie.
Pozwala to poprawnie renderować szeregi czasowe i mapę cieplną także wtedy, gdy
panel uruchamiany jest lokalnie na pobranej (i tym samym lekko nieaktualnej)
kopii bazy danych.

\subsection{Wspólne rozwiązywanie ścieżki bazy}

Funkcja \texttt{\_resolve\_db\_path()} kolejno sprawdza: zmienną środowiskową
\texttt{HONEYWATCH\_DB}, ścieżkę produkcyjną \texttt{/opt/honeywatch/honeywatch.db},
katalog projektu oraz pulpit użytkownika. Dzięki temu \textbf{ten sam kod}
uruchamia się bez zmian zarówno na serwerze produkcyjnym, jak i na komputerze
deweloperskim z lokalną kopią bazy.

\section{Wdrożenie produkcyjne}

System działa na serwerze VPS z systemem Linux. Panel jest utrzymywany przez
usługę \texttt{systemd} (\texttt{honeywatch-dashboard.service}), która
automatycznie wznawia jego działanie po awarii. Parser uruchamiany jest co
minutę przez \texttt{cron}. Skrypt \texttt{install-vps.sh} automatyzuje
instalację: instaluje zależności systemowe, tworzy środowisko wirtualne
Pythona, instaluje pakiety z \texttt{requirements.txt} oraz rejestruje
i uruchamia usługę. Listing \ref{lst:svc} przedstawia definicję usługi systemd.

\begin{lstlisting}[caption={Definicja usługi systemd panelu.},label={lst:svc}]
[Service]
Type=simple
WorkingDirectory=/opt/honeywatch/HoneyWatchApp
Environment=HONEYWATCH_DB=/opt/honeywatch/honeywatch.db
Environment=FLASK_DEBUG=0
Environment=PORT=5000
ExecStart=/opt/honeywatch/venv/bin/python -m dashboard.app
Restart=on-failure
RestartSec=5
WantedBy=multi-user.target
\end{lstlisting}

Ze względów bezpieczeństwa port panelu (5000) \textbf{nie jest} otwierany dla
całego świata --- baza zawiera adresy IP napastników oraz przechwycone hasła.
Dostęp do panelu ograniczono regułą zapory (\texttt{ufw}) do zaufanego adresu
administratora.

% =====================================================================
%  ROZDZIAŁ 4 — FUNKCJONALNOŚĆ Z PERSPEKTYWY UŻYTKOWNIKA
% =====================================================================
\chapter{Funkcjonalność z perspektywy użytkownika}
\label{ch:func}

Cały panel HoneyWatch został zaprojektowany jako \textbf{pojedyncza, przewijana
strona} (ang. \emph{single-page dashboard}) z ciemnym motywem kolorystycznym
i akcentem w barwie miodu (\textcolor{honey}{$\blacksquare$} \texttt{\#F59E0B}),
nawiązującym do nazwy projektu. Po lewej stronie znajduje się stały pasek
nawigacyjny, którego pozycje przewijają widok do odpowiednich sekcji. U góry
umieszczono pasek tytułowy z przełącznikiem zakresu czasu oraz przyciskiem
odświeżania. Rysunek \ref{fig:overview} przedstawia ogólny widok panelu.

\begin{figure}[ht]
\centering
\zrzut{8cm}{Ogólny widok panelu HoneyWatch (pasek boczny, KPI, mapa, wykresy).}
\caption{Pełny widok panelu głównego HoneyWatch.}
\label{fig:overview}
\end{figure}

\textit{Uwaga: w miejscach oznaczonych jako ,,Miejsce na zrzut ekranu''
należy wstawić rzeczywiste zrzuty z działającej aplikacji. Aby to zrobić, wystarczy
zastąpić polecenie \texttt{\textbackslash zrzut\{...\}\{...\}} poleceniem
\texttt{\textbackslash includegraphics[width=\textbackslash textwidth]\{screenshots/plik.png\}}.}

\section{Pasek nawigacyjny i nagłówek}

Pasek boczny zawiera logo projektu (sześciokąt nawiązujący do plastra miodu)
oraz pozycje nawigacyjne: \emph{Overview}, \emph{Attacks}, \emph{Map},
\emph{Statistics}, \emph{Logs} oraz \emph{Top Credentials}. Kliknięcie pozycji
płynnie przewija stronę do odpowiedniej sekcji, a aktywna sekcja jest
automatycznie podświetlana podczas przewijania (mechanizm
\texttt{IntersectionObserver}).

Nagłówek mieści \textbf{przełącznik zakresu czasu} z czterema opcjami:
\emph{1 godzina}, \emph{24h}, \emph{Tydzień} oraz \emph{Miesiąc}. Wybór zakresu
przeładowuje stronę z odpowiednim parametrem, a wszystkie wskaźniki, wykresy
i tabele dostosowują się do wybranego okna czasowego.

\begin{figure}[ht]
\centering
\zrzut{3cm}{Nagłówek z przełącznikiem zakresu czasu (1h / 24h / Tydzień / Miesiąc) i przyciskiem odświeżania.}
\caption{Pasek nagłówka z przełącznikiem zakresu czasu.}
\label{fig:header}
\end{figure}

\section{Wskaźniki KPI}

Na samej górze treści wyświetlane są cztery kafelki ze wskaźnikami kluczowymi
(KPI), dające natychmiastowy obraz sytuacji:

\begin{itemize}
  \item \textbf{Total Attacks} --- łączna liczba zarejestrowanych zdarzeń
        w wybranym oknie czasu;
  \item \textbf{Unique IPs} --- liczba unikalnych adresów źródłowych;
  \item \textbf{Login Attempts} --- liczba prób logowania (udanych i nieudanych);
  \item \textbf{Countries} --- liczba krajów, z których pochodziły ataki.
\end{itemize}

Każdy kafelek zawiera dodatkowo \textbf{wskaźnik zmiany procentowej} względem
poprzedniego, analogicznego okna czasu --- ze strzałką w górę (kolor czerwony,
wzrost zagrożenia) lub w dół (kolor zielony, spadek). Dzięki temu użytkownik od
razu widzi, czy aktywność rośnie, czy maleje.

\begin{figure}[ht]
\centering
\zrzut{3.5cm}{Cztery kafelki KPI: Total Attacks, Unique IPs, Login Attempts, Countries --- każdy ze wskaźnikiem zmiany \%.}
\caption{Panel wskaźników KPI.}
\label{fig:kpi}
\end{figure}

\section{Interaktywna mapa ataków}

Centralnym i najbardziej charakterystycznym elementem panelu jest interaktywna
mapa świata (Leaflet) z ciemnymi kafelkami CARTO. Każdy kraj będący źródłem
ataków jest oznaczony \textbf{okrągłym znacznikiem}, którego promień jest
proporcjonalny do liczby ataków --- im większe koło, tym intensywniejszy ruch
z danego regionu. Po najechaniu kursorem na znacznik wyświetlana jest nazwa
kraju oraz dokładna liczba zdarzeń, a panel informacyjny obok mapy aktualizuje
swoją treść. Mapa automatycznie dopasowuje widok tak, aby objąć wszystkie
znaczniki.

Obok mapy znajduje się lista \textbf{Top Countries} --- ranking krajów według
liczby ataków, z flagami i wartościami liczbowymi, uzupełniony pozycją zbiorczą
,,Other'' (Pozostałe).

\begin{figure}[ht]
\centering
\zrzut{8cm}{Interaktywna mapa świata ze znacznikami proporcjonalnymi do liczby ataków oraz lista Top Countries.}
\caption{Mapa źródeł ataków wraz z rankingiem krajów.}
\label{fig:map}
\end{figure}

\section{Wykresy i statystyki}

Sekcja statystyk prezentuje dane w formie czterech wykresów (Chart.js).

\subsection{Oś czasu ataków}

Wykres liniowy z gradientowym wypełnieniem pokazuje liczbę ataków w czasie.
Ziarnistość osi dostosowuje się do wybranego zakresu: minuty (1 godzina),
godziny (24h) lub dni (tydzień, miesiąc). Brakujące przedziały są wypełniane
zerami, dzięki czemu wykres jest ciągły i czytelny.

\begin{figure}[ht]
\centering
\zrzut{5cm}{Wykres liniowy ,,Attacks Over Time'' z gradientowym wypełnieniem.}
\caption{Oś czasu aktywności ataków.}
\label{fig:timeline}
\end{figure}

\subsection{Najczęstsze nazwy użytkowników i hasła}

Dwa poziome wykresy słupkowe prezentują rankingi najczęściej próbowanych
\textbf{nazw użytkowników} (kolor miodowy) oraz \textbf{haseł} (kolor czerwony).
Są to jedne z najciekawszych danych w całym systemie --- pokazują, jakich
poświadczeń ,,na ślepo'' próbują boty (klasycznie: \texttt{root}, \texttt{admin},
\texttt{123456}, \texttt{password}).

\begin{figure}[ht]
\centering
\zrzut{5cm}{Dwa poziome wykresy słupkowe: Top Usernames oraz Top Passwords.}
\caption{Rankingi najczęściej próbowanych poświadczeń.}
\label{fig:creds}
\end{figure}

\subsection{Typy zdarzeń}

Wykres pierścieniowy (\emph{doughnut}) przedstawia rozkład typów zdarzeń
(np. skanowanie portów, nieudane logowanie, udane logowanie, sondowanie HTTP),
z czytelną legendą HTML zawierającą udziały procentowe. Każdy typ ma przypisany
stały kolor, spójny w całym panelu.

\begin{figure}[ht]
\centering
\zrzut{5cm}{Wykres pierścieniowy ,,Event Types'' z legendą i udziałami procentowymi.}
\caption{Rozkład typów zdarzeń.}
\label{fig:events}
\end{figure}

\subsection{Mapa cieplna aktywności}

Mapa cieplna (\emph{heatmap}) w układzie 7 dni $\times$ 24 godziny pokazuje,
w których dniach tygodnia i o których godzinach aktywność jest największa.
Intensywność koloru każdej komórki odpowiada liczbie zdarzeń. Pozwala to wykryć
wzorce czasowe --- np. czy boty działają równomiernie przez całą dobę, czy
występują wyraźne szczyty.

\begin{figure}[ht]
\centering
\zrzut{4.5cm}{Mapa cieplna aktywności w układzie dzień tygodnia $\times$ godzina.}
\caption{Mapa cieplna aktywności ataków.}
\label{fig:heatmap}
\end{figure}

\section{Ostatnie ataki}

Sekcja ,,Recent Attacks'' wyświetla listę najnowszych zdarzeń w przystępnej
formie: flaga kraju, adres IP, typ zdarzenia oraz czas. Pozwala to na bieżąco
obserwować, co dzieje się ,,tu i teraz''.

\begin{figure}[ht]
\centering
\zrzut{4cm}{Lista ostatnich ataków z flagami krajów, adresami IP i typami zdarzeń.}
\caption{Lista ostatnich ataków.}
\label{fig:recent}
\end{figure}

\section{Przeszukiwalna tabela logów}

Na dole panelu znajduje się pełna \textbf{tabela logów} (\emph{Attack Logs})
prezentująca surowe zdarzenia w formie tabelarycznej, z kolumnami: czas, adres
IP, port, kraj (z flagą), typ zdarzenia, nazwa użytkownika, hasło oraz źródło
(honeypot). Tabela jest \textbf{stronicowana} (domyślnie 100 wierszy na stronę)
z przyciskami ,,Poprzednia''/,,Następna'' oraz licznikiem ,,od--do z N''.
Wszystkie wartości pochodzące od napastnika są bezpiecznie ,,ucieczkowane''
(\texttt{escapeHtml}), co zapobiega atakom XSS poprzez spreparowane dane
wejściowe.

\begin{figure}[ht]
\centering
\zrzut{7cm}{Stronicowana tabela logów z kolumnami czas/IP/port/kraj/typ/login/hasło/źródło.}
\caption{Przeszukiwalna, stronicowana tabela logów.}
\label{fig:logs}
\end{figure}

% =====================================================================
%  ROZDZIAŁ 5 — PRZYKŁADY UŻYCIA I DEMONSTRACJA
% =====================================================================
\chapter{Przykłady użycia i demonstracja wyników}
\label{ch:demo}

W tym rozdziale przedstawiamy kilka reprezentatywnych scenariuszy korzystania
z systemu wraz z uzyskanymi wynikami, które potwierdzają, że zaimplementowane
funkcje rzeczywiście działają na rzeczywistych danych zebranych przez honeypoty.

\section{Charakterystyka zebranych danych}

System był wystawiony na publiczny adres IP i przez okres działania zarejestrował
ataki z wielu krajów świata. Pamięć podręczna geolokalizacji
(\texttt{geo\_cache.json}) zawiera adresy IP m.in. z następujących krajów:
Holandia (NL), Stany Zjednoczone (US), Chiny (CN), Tajlandia (TH), Niemcy (DE),
Wielka Brytania (GB), Rosja (RU), Japonia (JP), Francja (FR), Brazylia (BR),
Polska (PL) oraz Iran (IR). Już ten zbiór dobrze ilustruje globalny
i zautomatyzowany charakter zjawiska --- ataki nadchodzą z całego świata,
przez całą dobę, bez żadnej legalnej przyczyny do łączenia się z naszym serwerem.

\section{Scenariusz 1: Szybka ocena sytuacji (KPI + mapa)}

\textbf{Cel:} analityk chce w kilka sekund ocenić bieżącą skalę zagrożenia.

\textbf{Kroki:}
\begin{enumerate}
  \item Otwarcie panelu w przeglądarce.
  \item Odczyt czterech wskaźników KPI (łączna liczba ataków, unikalne IP,
        próby logowania, liczba krajów) wraz ze wskaźnikami zmiany procentowej.
  \item Spojrzenie na mapę --- największe koła wskazują kraje generujące
        najwięcej ruchu.
  \item (Opcjonalnie) zmiana zakresu czasu na ,,1 godzina'', aby zobaczyć
        wyłącznie najświeższą aktywność.
\end{enumerate}

\textbf{Wynik:} natychmiastowy, syntetyczny obraz sytuacji bez konieczności
analizowania surowych logów. Wzrost wskaźnika ,,Total Attacks'' o wartość dodatnią
(czerwona strzałka) sygnalizuje nasilenie aktywności.

\section{Scenariusz 2: Analiza słownika botów (Top Credentials)}

\textbf{Cel:} sprawdzenie, jakich nazw użytkowników i haseł najczęściej próbują
napastnicy --- bezpośrednio przekłada się to na rekomendacje dotyczące polityki
haseł.

\textbf{Kroki:}
\begin{enumerate}
  \item Przewinięcie do sekcji wykresów słupkowych \emph{Top Usernames}
        i \emph{Top Passwords} (lub kliknięcie ,,Top Credentials'' w nawigacji).
  \item Odczyt rankingu --- typowo dominują wartości takie jak \texttt{root},
        \texttt{admin}, \texttt{user}, \texttt{123456}, \texttt{password},
        \texttt{12345}.
\end{enumerate}

\textbf{Wynik:} potwierdzenie, że ataki słownikowe opierają się na niewielkim
zestawie domyślnych i banalnych poświadczeń. \textbf{Wniosek obronny:} samo
wyłączenie logowania na konto \texttt{root} oraz wymuszenie silnych haseł
neutralizuje zdecydowaną większość obserwowanych prób.

\section{Scenariusz 3: Wykrywanie wzorców czasowych (heatmap + oś czasu)}

\textbf{Cel:} ustalenie, czy aktywność ataków ma charakter ciągły, czy
występują wyraźne szczyty.

\textbf{Kroki:}
\begin{enumerate}
  \item Ustawienie zakresu czasu na ,,Tydzień'' lub ,,Miesiąc''.
  \item Analiza wykresu osi czasu (trend ogólny) oraz mapy cieplnej
        (rozkład dzień $\times$ godzina).
\end{enumerate}

\textbf{Wynik:} zazwyczaj obserwuje się dość równomierny rozkład w ciągu doby,
co potwierdza zautomatyzowany (botowy), a nie ręczny charakter większości
ataków --- maszyny nie ,,śpią''.

\section{Scenariusz 4: Dochodzenie dotyczące pojedynczego adresu IP (tabela logów)}

\textbf{Cel:} szczegółowe prześledzenie aktywności wybranego adresu IP.

\textbf{Kroki:}
\begin{enumerate}
  \item Przewinięcie do sekcji \emph{Attack Logs}.
  \item Przeglądanie kolejnych stron tabeli (przyciski ,,Następna''/,,Poprzednia'').
  \item Odczyt szczegółów zdarzeń: portu, typu, próbowanych poświadczeń
        i honeypotu, który zarejestrował zdarzenie.
\end{enumerate}

\textbf{Wynik:} pełny, surowy materiał dowodowy dla każdego zdarzenia --- przydatny
np. do zgłoszenia nadużycia do dostawcy danej sieci lub do dodania adresu na
listę blokowanych.

\section{Przykładowe zapytania analityczne (SQL)}

Ponieważ dane są przechowywane w relacyjnej bazie SQLite, można je również
analizować bezpośrednio zapytaniami SQL. Poniższe przykłady odpowiadają
agregacjom realizowanym wewnętrznie przez panel.

\begin{lstlisting}[style=sqlstyle,caption={Najczęściej próbowane hasła.}]
SELECT password, COUNT(*) AS cnt
FROM attacks
WHERE password IS NOT NULL AND password <> ''
GROUP BY password
ORDER BY cnt DESC
LIMIT 10;
\end{lstlisting}

\begin{lstlisting}[style=sqlstyle,caption={Liczba unikalnych adresów IP wg kraju (po wzbogaceniu geolokalizacją).}]
SELECT src_ip, COUNT(*) AS cnt
FROM attacks
GROUP BY src_ip
ORDER BY cnt DESC;
\end{lstlisting}

\begin{lstlisting}[style=sqlstyle,caption={Aktywność w rozbiciu na godziny doby.}]
SELECT strftime('%H', timestamp) AS hour, COUNT(*) AS events
FROM attacks
GROUP BY hour
ORDER BY hour;
\end{lstlisting}

\section{Przykładowy wpis w logu honeypotu}

Listing \ref{lst:cowrie} pokazuje przykładowy wpis JSON z honeypotu Cowrie
(nieudana próba logowania), który parser przekształca w wiersz tabeli
\texttt{attacks} (\texttt{event\_type = login\_failed}, \texttt{success = 0}).

\begin{lstlisting}[caption={Przykładowy wpis JSON z logu Cowrie.},label={lst:cowrie}]
{"eventid": "cowrie.login.failed",
 "username": "root", "password": "123456",
 "src_ip": "185.224.128.17", "src_port": 54422,
 "timestamp": "2026-05-30T22:14:07.512Z",
 "session": "a1b2c3d4"}
\end{lstlisting}

% =====================================================================
%  ROZDZIAŁ 6 — ORGANIZACJA PRACY ZESPOŁU
% =====================================================================
\chapter{Organizacja pracy zespołu}
\label{ch:team}

Projekt był realizowany przez dwuosobowy zespół. Pracę zorganizowano w oparciu
o system kontroli wersji \textbf{Git} oraz platformę \textbf{GitHub}, co
umożliwiło równoległą pracę nad różnymi częściami systemu i bieżący przegląd
wprowadzanych zmian.

\section{Podział zadań i obszarów odpowiedzialności}

\begin{table}[ht]
\centering
\small
\begin{tabularx}{\textwidth}{>{\raggedright\arraybackslash}p{3.6cm} X}
\toprule
\textbf{Członek zespołu} & \textbf{Główne obszary odpowiedzialności} \\
\midrule
\textbf{Mateusz Rup} &
Architektura systemu i potoku danych; konfiguracja i wdrożenie honeypotów
(Cowrie, OpenCanary) na serwerze VPS; parser logów (\texttt{parser.py})
z przyrostowym czytaniem i deduplikacją; projekt schematu bazy danych; warstwa
dostępu do danych i logika analityczna (\texttt{db.py}); moduł geolokalizacji
(\texttt{geoip.py}); aplikacja serwerowa Flask (\texttt{app.py}); automatyzacja
wdrożenia (systemd, cron, skrypt instalacyjny); dokumentacja. \\
\midrule
\textbf{Jakub Gwizdała} &
Warstwa prezentacji i wizualizacji danych; integracja interaktywnej mapy
\textbf{Leaflet} ze znacznikami proporcjonalnymi do liczby ataków; współpraca
przy stylizacji frontendu i komponentach wizualnych panelu; testy interfejsu
użytkownika. \\
\bottomrule
\end{tabularx}
\caption{Podział obowiązków w zespole.}
\label{tab:team}
\end{table}

Podział ten odzwierciedlają również historia repozytorium i przepływ pracy
oparty na gałęziach funkcyjnych --- m.in. mapa Leaflet powstała na osobnej
gałęzi \texttt{feature/leaflet-map} i została włączona do gałęzi głównej poprzez
\emph{Pull Request} po przeglądzie kodu.

\section{Metody komunikacji i koordynacji}

\begin{itemize}
  \item \textbf{Git/GitHub} --- wspólne repozytorium, gałęzie funkcyjne,
        \emph{Pull Requesty} z przeglądem kodu przed scaleniem do gałęzi głównej.
  \item \textbf{Bieżąca komunikacja} --- regularne ustalenia dotyczące zakresu
        prac, podziału zadań i integracji poszczególnych komponentów.
  \item \textbf{Wspólny serwer VPS} --- jedno środowisko produkcyjne, na którym
        testowano działanie całego potoku na rzeczywistych danych.
\end{itemize}

\section{Przyjęty proces pracy}

Prace przebiegały etapowo, zgodnie z architekturą potoku danych: najpierw
uruchomiono i skonfigurowano honeypoty oraz zaprojektowano schemat bazy,
następnie powstał parser logów, później warstwa dostępu do danych i aplikacja
serwerowa, a na końcu warstwa wizualizacji (mapa, wykresy) oraz testy całości
na żywych danych. Takie podejście pozwalało weryfikować poprawność każdego etapu
na rzeczywistych danych, zanim przystąpiono do kolejnego.

% =====================================================================
%  ROZDZIAŁ 7 — PODSUMOWANIE I WNIOSKI
% =====================================================================
\chapter{Podsumowanie i wnioski}
\label{ch:summary}

\section{Osiągnięte rezultaty}

Wszystkie założone cele projektu zostały zrealizowane. Powstał kompletny,
działający system, który:

\begin{itemize}
  \item zbiera \textbf{rzeczywiste} dane o atakach z publicznie wystawionego
        serwera, wykorzystując dwa uzupełniające się honeypoty (Cowrie
        i OpenCanary);
  \item przetwarza surowe logi w ustrukturyzowaną bazę danych za pomocą
        samodzielnie napisanego, odpornego na błędy i duplikaty parsera;
  \item wzbogaca dane o geolokalizację adresów IP z poszanowaniem limitów
        zewnętrznego API (trwała pamięć podręczna);
  \item prezentuje wyniki w czytelnym, interaktywnym panelu webowym
        (mapa, wykresy, mapa cieplna, stronicowana tabela logów);
  \item działa w trybie produkcyjnym, automatycznie aktualizując dane
        (cron + systemd).
\end{itemize}

\section{Wnioski techniczne}

\begin{itemize}
  \item \textbf{Prostota ma wartość.} Rezygnacja z zależności zewnętrznych
        w parserze oraz wybór lekkiej bazy SQLite znacznie uprościły wdrożenie
        i utrzymanie systemu na minimalnym serwerze VPS.
  \item \textbf{Idempotentność się opłaca.} Mechanizm przyrostowego czytania
        logów połączony z deduplikacją (\texttt{INSERT OR IGNORE} + unikalne
        indeksy) sprawia, że parser można uruchamiać dowolnie często, bez
        ryzyka zniekształcenia danych.
  \item \textbf{Bezpośredni odczyt bazy} przez panel okazał się trafną decyzją
        architektoniczną --- wyeliminował potrzebę budowania pośredniej warstwy
        API i synchronizacji danych między procesami.
\end{itemize}

\section{Wnioski merytoryczne (dotyczące zjawiska)}

Zebrane dane potwierdziły, że \textbf{każdy} serwer wystawiony do Internetu jest
niemal natychmiast i nieprzerwanie atakowany przez zautomatyzowane boty z całego
świata. Ataki opierają się głównie na prostych słownikach domyślnych poświadczeń,
co oznacza, że podstawowe środki higieny bezpieczeństwa (silne, unikalne hasła,
wyłączenie logowania na konto \texttt{root}, ograniczenie dostępu do usług
administracyjnych) neutralizują zdecydowaną większość obserwowanego ruchu.

\section{Ograniczenia rozwiązania}

W duchu rzetelności wskazujemy znane ograniczenia obecnej wersji:

\begin{itemize}
  \item \textbf{Brak automatycznego odświeżania w przeglądarce} --- aktualizacja
        danych wymaga odświeżenia strony (przyciskiem lub klawiszem F5).
  \item \textbf{Geolokalizacja zależna od zewnętrznego API} --- dokładność
        i dostępność danych o położeniu zależą od usługi ip-api.com.
  \item \textbf{Honeypoty średnio-interaktywne} --- system rejestruje próby
        i polecenia, lecz nie udostępnia napastnikowi pełnego środowiska, co
        ogranicza głębię analizy zaawansowanych ataków.
  \item \textbf{Brak warstwy uwierzytelniania panelu} --- dostęp jest
        zabezpieczony wyłącznie na poziomie zapory sieciowej (reguła \texttt{ufw}).
  \item \textbf{Skala} --- architektura jednoserwerowa oparta na SQLite jest
        odpowiednia dla tego projektu, lecz nie jest przeznaczona do obsługi
        wielu honeypotów rozproszonych geograficznie.
  \item \textbf{Brak zautomatyzowanych testów} --- weryfikacja odbywała się
        manualnie na rzeczywistych danych.
\end{itemize}

\section{Cechy szczególne i innowacyjne}

Na tle istniejących, ciężkich platform (np. T-Pot) HoneyWatch wyróżnia się
\textbf{lekkością i czytelnością}: cały potok przetwarzania danych został
napisany od podstaw, bez zbędnych zależności, a panel kładzie nacisk na
przejrzystą prezentację najważniejszych informacji. Mechanizmy takie jak
przyrostowe, idempotentne parsowanie, niezależność od formatu znacznika czasu
czy współdzielona ścieżka bazy (jeden kod dla produkcji i dla środowiska
deweloperskiego) stanowią konkretne, przemyślane rozwiązania inżynierskie.

\section{Kierunki dalszego rozwoju}

\begin{itemize}
  \item automatyczne odświeżanie danych w przeglądarce (np. WebSocket lub
        cykliczne odpytywanie API);
  \item dodanie kolejnych honeypotów oraz dokończenie obsługi Dionaea;
  \item filtrowanie i wyszukiwanie w tabeli logów (po IP, kraju, dacie);
  \item eksport danych i raportów (CSV, PDF);
  \item warstwa uwierzytelniania dostępu do panelu;
  \item wzbogacenie danych o reputację IP (np. listy znanych botnetów)
        oraz powiadomienia o nietypowych skokach aktywności.
\end{itemize}

% =====================================================================
%  BIBLIOGRAFIA (styl IEEE)
% =====================================================================
\begin{thebibliography}{9}
\bibitem{cowrie}
Cowrie Authors, \emph{Cowrie SSH/Telnet Honeypot --- Documentation}. [Online].
Dostępne: \url{https://cowrie.readthedocs.io}. [Dostęp: 8 cze. 2026].

\bibitem{opencanary}
Thinkst Applied Research, \emph{OpenCanary Documentation}. [Online].
Dostępne: \url{https://opencanary.readthedocs.io}. [Dostęp: 8 cze. 2026].

\bibitem{flask}
Pallets Projects, \emph{Flask Documentation (3.x)}. [Online].
Dostępne: \url{https://flask.palletsprojects.com}. [Dostęp: 8 cze. 2026].

\bibitem{sqlite}
SQLite Consortium, \emph{SQLite Documentation}. [Online].
Dostępne: \url{https://www.sqlite.org/docs.html}. [Dostęp: 8 cze. 2026].

\bibitem{python}
Python Software Foundation, \emph{The Python Standard Library}. [Online].
Dostępne: \url{https://docs.python.org/3/library/}. [Dostęp: 8 cze. 2026].

\bibitem{leaflet}
V. Agafonkin, \emph{Leaflet --- an open-source JavaScript library for
interactive maps}. [Online]. Dostępne: \url{https://leafletjs.com}.
[Dostęp: 8 cze. 2026].

\bibitem{chartjs}
Chart.js Contributors, \emph{Chart.js Documentation}. [Online].
Dostępne: \url{https://www.chartjs.org/docs/latest/}. [Dostęp: 8 cze. 2026].

\bibitem{ipapi}
ip-api.com, \emph{IP Geolocation API Documentation}. [Online].
Dostępne: \url{https://ip-api.com/docs}. [Dostęp: 8 cze. 2026].

\bibitem{carto}
CARTO, \emph{CARTO Basemaps (Dark Matter)}. [Online].
Dostępne: \url{https://carto.com/basemaps/}. [Dostęp: 8 cze. 2026].

\bibitem{spitzner}
L. Spitzner, \emph{Honeypots: Tracking Hackers}. Boston, MA, USA:
Addison-Wesley, 2002.
\end{thebibliography}

% =====================================================================
%  ZAŁĄCZNIK
% =====================================================================
\appendix
\chapter{Załącznik: kod źródłowy}
\label{ch:appendix}

Kompletny, działający i skomentowany kod źródłowy projektu jest dostępny
w publicznym repozytorium Git:

\begin{center}
\large\url{https://github.com/HoneyWatch/HoneyWatchApp}
\end{center}

\section{Struktura repozytorium}

\begin{lstlisting}[caption={Struktura katalogów projektu.}]
HoneyWatchApp/
|-- README.md                 # Opis projektu
|-- requirements.txt          # Zaleznosci (Flask)
|-- dashboard/                # Aplikacja webowa (Flask)
|   |-- app.py                # Endpointy: 1 strona HTML + API JSON
|   |-- db.py                 # Warstwa danych i logika analityczna
|   |-- geoip.py              # Geolokalizacja IP + cache
|   |-- iso_codes.py          # Mapowanie ISO alpha2 -> alpha3
|   |-- templates/            # Szablony Jinja2
|   `-- static/               # CSS + JavaScript (mapa, wykresy)
|-- data/
|   `-- geo_cache.json        # Pamiec podreczna geolokalizacji
`-- deploy/                   # Wdrozenie na VPS
    |-- parser.py             # Parser logow (Cowrie + OpenCanary)
    |-- install-vps.sh        # Skrypt instalacyjny
    |-- honeywatch-dashboard.service  # Usluga systemd
    `-- VPS.md                # Instrukcja wdrozeniowa
\end{lstlisting}

\section{Uruchomienie lokalne}

\begin{lstlisting}[style=bashstyle,caption={Uruchomienie panelu w środowisku deweloperskim.}]
# 1. Instalacja zaleznosci
python -m venv venv
venv\Scripts\activate          # Windows
pip install -r requirements.txt

# 2. Wskazanie kopii bazy danych (opcjonalnie)
set HONEYWATCH_DB=C:\sciezka\do\honeywatch.db

# 3. Uruchomienie panelu
python -m dashboard.app
# Panel dostepny pod http://localhost:5000
\end{lstlisting}

\section{Uruchomienie parsera (serwer VPS)}

\begin{lstlisting}[style=bashstyle,caption={Cykliczne uruchamianie parsera przez cron.}]
# Wpis w crontab (uruchomienie co minute)
* * * * * /opt/honeywatch/venv/bin/python /opt/honeywatch/parser.py
\end{lstlisting}

\end{document}






